% definitions.tex — Book I definitions (the 23). Each is a `scope`
% block: definitions delimit terms used by later claims and edges, but
% they're not themselves assertions about the world.

\section*{Book I: Definitions}
\label{sec:def-bookI}

\begin{scope}[Definition I.1: point]
\label{def:I.1}
A point is that which has no part.
\end{scope}

\begin{scope}[Definition I.2: line]
\label{def:I.2}
A line is breadthless length.
\end{scope}

\begin{scope}[Definition I.3: extremities of a line]
\label{def:I.3}
The extremities of a line are points.
\end{scope}

\begin{scope}[Definition I.4: straight line]
\label{def:I.4}
A straight line is a line which lies evenly with the points on itself.
\end{scope}

\begin{scope}[Definition I.5: surface]
\label{def:I.5}
A surface is that which has length and breadth only.
\end{scope}

\begin{scope}[Definition I.6: extremities of a surface]
\label{def:I.6}
The extremities of a surface are lines.
\end{scope}

\begin{scope}[Definition I.7: plane surface]
\label{def:I.7}
A plane surface is a surface which lies evenly with the straight lines
on itself.
\end{scope}

\begin{scope}[Definition I.8: plane angle]
\label{def:I.8}
A plane angle is the inclination to one another of two lines in a
plane which meet one another and do not lie in a straight line.
\end{scope}

\begin{scope}[Definition I.9: rectilineal angle]
\label{def:I.9}
And when the lines containing the angle are straight, the angle is
called rectilineal.
\end{scope}

\begin{scope}[Definition I.10: right angle]
\label{def:I.10}
When a straight line set up on a straight line makes the adjacent
angles equal to one another, each of the equal angles is right, and
the straight line standing on the other is called a perpendicular to
that on which it stands.
\end{scope}

\begin{scope}[Definition I.11: obtuse angle]
\label{def:I.11}
An obtuse angle is an angle greater than a right angle.
\end{scope}

\begin{scope}[Definition I.12: acute angle]
\label{def:I.12}
An acute angle is an angle less than a right angle.
\end{scope}

\begin{scope}[Definition I.13: boundary]
\label{def:I.13}
A boundary is that which is an extremity of anything.
\end{scope}

\begin{scope}[Definition I.14: figure]
\label{def:I.14}
A figure is that which is contained by any boundary or boundaries.
\end{scope}

\begin{scope}[Definition I.15: circle]
\label{def:I.15}
A circle is a plane figure contained by one line such that all the
straight lines falling upon it from one point among those lying within
the figure are equal to one another.
\end{scope}

\begin{scope}[Definition I.16: centre]
\label{def:I.16}
And the point is called the centre of the circle.
\end{scope}

\begin{scope}[Definition I.17: diameter]
\label{def:I.17}
A diameter of the circle is any straight line drawn through the centre
and terminated in both directions by the circumference of the circle,
and such a straight line also bisects the circle.
\end{scope}

\begin{scope}[Definition I.18: semicircle]
\label{def:I.18}
A semicircle is the figure contained by the diameter and the
circumference cut off by it. And the centre of the semicircle is the
same as that of the circle.
\end{scope}

\begin{scope}[Definition I.19: rectilineal figures]
\label{def:I.19}
Rectilineal figures are those which are contained by straight lines,
trilateral figures being those contained by three, quadrilateral those
contained by four, and multilateral those contained by more than four
straight lines.
\end{scope}

\begin{scope}[Definition I.20: kinds of trilateral]
\label{def:I.20}
Of trilateral figures, an equilateral triangle is that which has its
three sides equal, an isosceles triangle that which has two of its
sides alone equal, and a scalene triangle that which has its three
sides unequal.
\end{scope}

\begin{scope}[Definition I.21: kinds of trilateral by angle]
\label{def:I.21}
Further, of trilateral figures, a right-angled triangle is that which
has a right angle, an obtuse-angled triangle that which has an obtuse
angle, and an acute-angled triangle that which has its three angles
acute.
\end{scope}

\begin{scope}[Definition I.22: kinds of quadrilateral]
\label{def:I.22}
Of quadrilateral figures, a square is that which is both equilateral
and right-angled; an oblong that which is right-angled but not
equilateral; a rhombus that which is equilateral but not right-angled;
and a rhomboid that which has its opposite sides and angles equal to
one another but is neither equilateral nor right-angled. And let
quadrilaterals other than these be called trapezia.
\end{scope}

\begin{scope}[Definition I.23: parallel straight lines]
\label{def:I.23}
Parallel straight lines are straight lines which, being in the same
plane and being produced indefinitely in both directions, do not meet
one another in either direction.
\end{scope}
